Los resultados obtenidos en esta tarea nos permiten concluir que la programación dinámica no solo es teóricamente más eficiente que la fuerza bruta, sino que también se manifiesta como la opción superior al momento de resolver instancias prácticas del problema de diferencias entre secuencias. La incapacidad de la fuerza bruta para escalar más allá de entradas pequeñas valida las predicciones sobre su complejidad exponencial, mientras que la solución dinámica mantiene tiempos de ejecución bajos y consistentes incluso en los casos más grandes evaluados.

Estos hallazgos reflejan la importancia de aplicar paradigmas algorítmicos adecuados según la naturaleza del problema, y evidencian cómo una estrategia permite abordar situaciones que serían intratables de otro modo. El contraste entre ambos enfoques también refuerza el valor del análisis experimental como herramienta complementaria al análisis teórico, especialmente cuando se desea observar el comportamiento real de algoritmos en condiciones controladas.

En definitiva, el experimento confirma que el paradigma de programación dinámica ofrece una solución viable, eficiente y reproducible para comparar secuencias mediante la subsecuencia común más larga, y que elegir el enfoque correcto marca una diferencia tanto en rendimiento como en aplicabilidad.
